% !TEX root = ../main.tex
\chapter{Conclusion}
\label{ch:conclusion}

\section{Summary of Achievements}

This project set out to design, implement, and evaluate a RISC-V-attached
hardware accelerator for CSR sparse matrix--dense matrix multiplication on the
Terasic DE1-SoC, and to demonstrate a measurable speedup over software while
remaining reproducible at the register-transfer level. That aim was met. A
complete PicoRV32 SoC with a working SpMM accelerator, 64\,KB shared
dual-port BRAM, MMIO interface, and UART reporting was synthesised and
deployed; three recorded optimisation stages (baseline, DSP-enhanced, and
INT8 packed) were each verified for correctness on hardware; and a 12-case
benchmark harness with deterministic matrix generation, cycle-accurate timing,
and element-by-element correctness checking was developed to evaluate them.

The final INT8-packed implementation achieves a peak speedup of
$\mathbf{56.80\times}$ over the quantised PicoRV32 software reference and
$\mathbf{11.92\times}$ over an ARM Cortex-M0 baseline running the same
algorithm, with an average of $45.28\times$ and $9.79\times$ respectively
across the 12-case suite. All 12 cases passed correctness checks at every
stage.

\section{Significance}

The significance of this work lies not in the absolute numbers compared to
large research accelerators, but in what the staged measurement reveals about
the workload. The $1.32\times$ gain from DSP-based MAC versus the
$2.18\times$ gain from INT8 packing experimentally confirms that the dominant
bottleneck in CSR SpMM on a memory-bandwidth-constrained embedded FPGA is
operand fetch traffic, not arithmetic execution. This is a useful result
beyond the specific artefact: it indicates that similar embedded accelerators
should prioritise format-level and precision-level optimisations over wider
MAC arrays. The project also demonstrates that a fully integrated, firmware-
driven hardware benchmark harness is a practical alternative to external
simulation infrastructure for small-scale accelerator evaluation, and that a
useful SpMM accelerator fits comfortably within the resource budget of a
Cyclone~V device.

\section{Future Work}

The most impactful extensions to the current design would be:

\begin{enumerate}
  \item \textbf{Wider tile width.} INT8 packing reduces the per-nonzero fetch
        cost enough that widening from $TN = 8$ to $TN = 16$ or $TN = 32$ is
        now attractive. Doubling the tile width costs an additional two
        32-bit reads per nonzero but delivers $2\times$ the MAC throughput
        per nonzero, and should push the speedup ceiling further.

  \item \textbf{Ping-pong $B$-segment prefetch.} Double-buffering the
        $B$-segment allows the next segment to be fetched while the current
        MAC operation completes, overlapping memory latency with arithmetic
        and removing it from the critical path.

  \item \textbf{External memory interface.} Extending the design to address
        off-chip SDRAM or HPS DRAM through the DE1-SoC's HPS bridge would
        permit realistically sized GNN and recommender-system matrices to
        be benchmarked, at which point off-chip bandwidth becomes the
        dominant term and different optimisations become relevant.

  \item \textbf{Interrupt-driven completion.} Replacing the polling loop
        with a RISC-V interrupt handler would free the CPU during
        accelerator execution, enabling pipelined benchmark preparation or
        post-processing of the previous result.

  \item \textbf{Power and energy characterisation.} Instrumenting the
        DE1-SoC rails or using Quartus Power Analyser with simulation-derived
        toggle files would add energy efficiency to the evaluation, which is
        the metric that matters most in the embedded and edge contexts for
        which this class of accelerator is targeted.
\end{enumerate}

\section{Closing Remarks}

This project has delivered a complete, measured, and reproducible RISC-V
FPGA-attached SpMM accelerator that achieves a peak speedup of $56.80\times$
on hardware through three progressively optimised designs. The work confirms
that even without exotic hardware resources or advanced scheduling, careful
alignment of the memory format, datapath structure, and quantisation strategy
with the actual bottleneck of the workload produces large, predictable, and
experimentally confirmed performance gains.
